% !TeX root = ../papers/nc-p3.tex
\section{反対階数14の二楕円成分の全通常点でのAS正則性}
\label{sec:rank-fourteen-all-ordinary}
\begin{theorem}\label{thm:rank-fourteen-all-ordinary}
標数零で定理\ref{thm:fixed-opposite-rank-drop}の二つの階数14成分に属する
全ての通常窓は、標準Koszul左右AS正則である。各始点の道空間は四次の中心元四本が
生成する多項式環上で階数256の自由加群である。特に、同定理の倍率方程式を満たす
各通常倍率の上の四つの階数14窓は全て正則である。
\end{theorem}
\begin{proof}
定理\ref{thm:rank-fourteen-ramified-disks}により $t$ が101で19または80に還元される
場合を扱えばよい。残る二値は $t^2+2t+5$ の相異なる根である。
$J^2=-1$、$\alpha=-1+2J$ と置き、$J$ の二つの101進埋込みを取る。
係数環を $D=\mathbb Z_{(101)}[J]$、局所変数を $u=t-\alpha$ とする。
成分の定数 $i$ には $\pm J$ を、$s$ の境界値には $\pm1$ を独立に取る。
$Y^2=2(t-1)(t+1)(t+3)$ の各解は、定数項
$Y_0=\pm(\alpha+1)(\alpha+3)$ から $D[[u]]$ に一意に持ち上がる。
実際 $2Y_0$ は単元である。通常な持上げでは $u$ の付値は正で、$u\ne0$ である。

定理\ref{thm:rank-fourteen-central-cut}で使った二次関係と四中心式を取る。
分母を払う際は中心元の全四位置で共通の倍率を取る。得られる多項式切断は中心である。
実際、長さ5の道空間は通常域で定階数であり、一般点で零の交換子は各点でも零である。
これらの切断を $t=\alpha+u$ に代入する。各始点・長さ・文字で前節までと同じ
右への一文字追加を行い、列に対応する語を定数項が単元のピボットで消す。
両複素埋込みを同時に扱うため、定数項のガウスノルムが101で単元かを検査する。

ピボットが見つからないときが境界の階数低下に対応する。その際は残る全行の
\emph{定数項が $\mathbb Q(J)$ 上で零}であることを直接確認し、各行を $u$ で一度割る。
これは $D[[u]][1/u]$ 上で有効な関係式である。各段の自由列が目標数まで減ったら
後退代入し、全列から残存語への射影を作る。一部の行を捨てても残存語は商を張る。
中心式の挿入は長さ4で行い、その後は矢との交換で全ての挿入を含む。

\path{computations/verify_rank_fourteen_all_ordinary.py} はこれを有限の厳密計算で行う。
初めに係数を $u^{27}$ を法として求める。単元の逆は同じ精度で計算でき、
$u$ による除算一回だけがその行を一桁短くする。ある文字ブロックで $k$ 回の
一斉除算を行うと、後退代入した全列も少なくとも元の精度から $k$ を引いた精度を持つ。
次の長さでは四ブロックの下限の最小値を使う。各除算時の零判定と単元判定には
定数項が完全に決まる精度が残る。このため、打切り展開から無限級数への外挿を要しない。
全16組（$i$、$s$ の符号と四始点）で長さ13まで正の精度が残り、商次元の上界は
\[
 1,4,10,20,31,40,44,40,31,20,10,4,1,0
\]
となる。各証明書の行操作は完備付値環で $u$ にその値を代入でき、必要な $u$ の逆だけを
分数体で取る。従って還元が19または80の通常窓にもこの上界が成り立つ。

未切断の道空間の各次数には一般正則点のHilbert次元が下界を与える。
四中心商から持ち上げた自由加群の上界がこの下界に一致し、
定理\ref{thm:klein-rational-period-four}の左右Koszul複体の証明が適用できる。
数体上の幾何点は101上の付値に埋め込み、体拡大による完全性を降下する。
超越的な $t$ はすでに正則な幾何学的一般点からの体拡大である。
\end{proof}
